%============================================================= SECTION 7: RESULTS
\section{Results, Success Analysis \& Empirical Verification Matrix}\label{sec:results}

This section answers requirements~(b) and~(c): whether each phase of the attack succeeded, the
technical reasons underlying successes and mitigations, what was observed across machines, and a
consolidated verification matrix across all laboratory runs.

\subsection{Per-Phase Outcome Verdict}
Table~\ref{tab:results_verdict} presents the definitive outcome of each functional reconnaissance
and exploitation module evaluated across the testbed:

\begin{table}[!ht]
\centering
\footnotesize
\renewcommand{\arraystretch}{1.06}
\caption{Systematic Operational Verdict and Evidence Across Attack Phases}
\label{tab:results_verdict}
\begin{tabularx}{\textwidth}{@{}L{3.6cm} L{2.6cm} X@{}}
\toprule
\hdrrow \thh{Reconnaissance / Exploit Phase} & \thh{Observed Outcome} & \thh{Empirical Ground-Truth Evidence}\\
\midrule
Phase 1: TCP SYN Port Scan & \textbf{Success} & Half-open classifications perfectly matched victim host \texttt{netstat} listeners; 100-port scale test correctly classified all ports.\\
\addlinespace[2pt]
Phase 2: Banner Grabbing & \textbf{Success} & Captured authentic daemon banners across FTP, HTTP, and Microsoft/Debian OpenSSH listeners.\\
\addlinespace[2pt]
Phase 3: OS Fingerprinting & \textbf{Success$^\ast$} & Correct OS family identified across all five testbed platforms; confidence reported honestly at 48\%--76\% ($^\ast$probabilistic).\\
\addlinespace[2pt]
Exploit 1: FTP Anonymous Login & \textbf{Success} & Authenticated with return code 230; successfully enumerated directories via passive mode \texttt{LIST}.\\
\addlinespace[2pt]
Exploit 2: FTP Web-Shell CGI RCE & \textbf{Success} & Uploaded \texttt{shell.py}; executed \texttt{whoami} returning \texttt{homed01\textbackslash turjo}; cleaned up via \texttt{DELE}.\\
\addlinespace[2pt]
Exploit 3: FTP Bounce Attack & \textbf{Mitigated} & RFC~959 \texttt{PORT} pivot aborted with \texttt{[WinError 10053]} due to host Windows Filtering Platform (WFP) socket policies.\\
\addlinespace[2pt]
Exploit 4: HTTP Directory Traversal & \textbf{Success} & Exfiltrated \texttt{win.ini} on Windows and \texttt{/etc/passwd} on Linux across standard, path-separated, and encoded vectors.\\
\addlinespace[2pt]
Exploit 5: SSH Weak Credentials & \textbf{Success} & Cracked \texttt{test:test123} on attempt 15/20; executed silent reconnaissance in quiet mode and multi-modal desktop payloads in loud mode.\\
\addlinespace[2pt]
Countermeasure: Active Defense & \textbf{Success (Bonus)} & Autonomous quarantine injected firewall rules dynamically, severing active attack chains mid-execution on Windows and Linux.\\
\bottomrule
\end{tabularx}
\end{table}

\subsection{Why the Attacks Succeeded}
Three engineering foundations enabled the toolkit to operate effectively:
\begin{enumerate}
    \item \textbf{High-Fidelity From-Scratch Framing:} Because all Ethernet, IPv4, and TCP headers
    are packed and parsed directly by our code, incoming frames are analyzed without normalization
    or masking by third-party abstractions.
    \item \textbf{Multi-Vector Probabilistic Synthesis:} By combining seven independent stack
    behaviors, the engine remains robust against network anomalies. If a firewall suppresses
    timestamps or normalizes TTLs, orthogonal signals (such as IP-ID sequencing or option ordering)
    maintain accurate attribution.
    \item \textbf{Dynamic OS-Aware Weaponization Loop:} Closing the loop between reconnaissance and
    exploitation allowed the identical tool binary to exfiltrate Windows configuration files on
    one host and Linux user credentials on another without manual intervention.
\end{enumerate}

\subsection{Analytical Assessment of Attack Observations \& Exploit Divergence}
In accordance with laboratory evaluation guidelines, we critically analyzed each attack vector
across five systematic criteria: (1)~\textbf{Attack Process}, (2)~\textbf{Assumptions Regarding the
Attack}, (3)~\textbf{Expected Outcome}, (4)~\textbf{Actual Observed Outcome}, and (5)~\textbf{Divergence
Analysis (Why It Differs)}.

\subsubsection{Deep-Dive on Exploit 3: RFC~959 FTP Bounce Attack (The Sole Exploit Failure)}
Across our exploitation battery, the \textbf{RFC~959 FTP Bounce Attack was the only module that did not succeed as an offensive exploit}. The exploit attempted to use the FTP daemon as an arbitrary TCP port scanner and proxy by establishing a control channel, issuing an RFC~959 \texttt{PORT 192,168,0,116,0,80} command toward an internal service, and triggering data transfer via \texttt{LIST}. Rather than completing a relayed TCP 3-way handshake on behalf of the attacker, the attack aborted with a transport reset:
\begin{center}\vspace{-3pt}
\small\texttt{500 Connection error: [WinError 10053] An established connection was aborted by the software in your host machine.}
\vspace{-3pt}\end{center}
This divergence highlights modern kernel-level stateful socket isolation. Specifically, the \textbf{Windows Filtering Platform (WFP)} and Winsock kernel drivers track session ownership across interface boundaries. When the FTP server process attempted an unauthorized cross-socket redirection to a third-party endpoint, the WFP filter detected the policy violation and issued an abortive socket reset (\texttt{WSAECONNABORTED}) before any outbound SYN could leave the network adapter. In hardened production environments, complementary application-layer RFC~2577 validation similarly ensures that daemons reject third-party \texttt{PORT} arguments when \texttt{sock.getpeername()} does not match the control endpoint.

\subsubsection{Comparative Observations Across All Attack Vectors}
Table~\ref{tab:attack_observations} provides the systematic comparative analysis across all evaluated modules according to the evaluation framework (Attack Process, Underlying Assumptions, Expected vs.\ Actual Outcomes, and Divergence Analysis):

\begin{table}[!ht]
\centering
\scriptsize
\renewcommand{\arraystretch}{1.06}
\caption{Systematic Attack Observations: Process, Assumptions, Expected vs.\ Actual Outcomes, and Divergence Analysis}
\label{tab:attack_observations}
\begin{tabularx}{\textwidth}{@{}L{2.8cm} X X@{}}
\toprule
\hdrrow \thh{Attack Vector} & \thh{Process \& Underlying Assumptions} & \thh{Expected vs.\ Actual Outcome \& Divergence Analysis}\\
\midrule
\textbf{FTP Anonymous Access} (Section~\ref{sec:exploits}) &
Submits \texttt{USER anonymous}; assumes server permits unauthenticated guest sessions and directory listing. &
\textbf{Expected:} Return 230, list directory.\newline
\textbf{Actual:} \textbf{Success} (Code 230; listed \texttt{index.html}, \texttt{pub/}, \texttt{uploads/}).\\
\addlinespace[2pt]
\textbf{FTP Web-Shell CGI RCE} (Section~\ref{sec:exploits}) &
Uploads \texttt{shell.py} via \texttt{STOR}, triggers \texttt{whoami} via HTTP GET; assumes writeable directory is served as CGI. &
\textbf{Expected:} Remote code execution.\newline
\textbf{Actual:} \textbf{Success} (Executed \texttt{whoami} $\to$ \texttt{homed01\textbackslash turjo}; payload deleted via \texttt{DELE}).\\
\addlinespace[2pt]
\textbf{RFC 959 FTP Bounce Attack} (Section~\ref{sec:exploits}) &
Sends \texttt{PORT 192,168,0,116,0,80} + \texttt{LIST}; assumes daemon relays data connections to third-party hosts without kernel isolation. &
\textbf{Expected:} Proxy port scan.\newline
\textbf{Actual:} \textbf{Failed / Mitigated} (\texttt{WinError 10053}; aborted by host WFP).\newline
\textbf{Why it differs:} Host WFP blocks socket relay.\\
\addlinespace[2pt]
\textbf{HTTP Directory Traversal} (Section~\ref{sec:exploits}) &
Issues relative dot-dot-slash paths; assumes web daemon lacks input path canonicalization and has OS file read permissions. &
\textbf{Expected:} Exfiltrate target config.\newline
\textbf{Actual:} \textbf{Success} (Retrieved \texttt{win.ini} on Win 11; \texttt{/etc/passwd} on Linux across 3 encodings).\\
\addlinespace[2pt]
\textbf{SSH Weak Credentials} (Section~\ref{sec:exploits}) &
Dictionary brute force over port~22; assumes password authentication allowed, weak password in dictionary, no rate limiting. &
\textbf{Expected:} Remote terminal access.\newline
\textbf{Actual:} \textbf{Success} (Cracked on attempt 15/20; deployed desktop artifacts \& voice alert).\\
\addlinespace[2pt]
\textbf{Active Defense Quarantining} (Section~\ref{sec:defense}) &
Passive L2 capture monitor, sliding-window rate tracking, canary traps; assumes attacker scans multiple ports rapidly. &
\textbf{Expected:} Auto-block attacker.\newline
\textbf{Actual:} \textbf{Success (Bonus)} (Severed attack chains at 8 SYNs on Win/Linux via dynamic kernel firewall rules).\\
\bottomrule
\end{tabularx}
\end{table}

\subsubsection{Probabilistic Confidence Calibration vs.\ False Certainty}
In our evaluation, OS fingerprinting confidence ranged from 48\% to 76\%, rather than reporting
an unrealistic 100\%. Initial TTL~64 is shared between Linux and Darwin, NAT gateways can normalize
IP IDs, and targets with closed ports offer sparse telemetry. The engine reports calibrated
probabilities and coarse OS tiers rather than claiming certainty where signals overlap.

\FloatBarrier
\subsection{Victim Server Daemon Observations}
The victim node hosted vulnerable lab services (\path{victim/http_server.py} on port~80 and
\path{victim/ftp_server.py} on port~21). The corresponding daemon logs below record the inbound
attacks from the server's perspective:

\begin{logbox}
[Victim Terminal 1: Vulnerable HTTP Server (Port 80)]
[HTTP] Listening on 0.0.0.0:80 | CGI enabled for uploads/*.py | Traversal protection DISABLED
[HTTP] ('192.168.0.105', 51482) -> HEAD /                               (Banner probe)
[HTTP] ('192.168.0.105', 51488) -> GET /uploads/shell.py?cmd=whoami      (CGI execution)
[HTTP] ('192.168.0.105', 51488) -> CGI executed shell.py (query='cmd=whoami') -> 200 OK
[HTTP] ('192.168.0.105', 51493) -> GET /..\..\..\..\..\windows\win.ini   -> 200 OK (Served)
[HTTP] ('192.168.0.105', 51498) -> GET /..%5c..%5c..%5cwindows%5cwin.ini -> 200 OK (Served)

[Victim Terminal 2: Vulnerable FTP Server (Port 21)]
[FTP] Listening on 0.0.0.0:21 | Root: ./victim/lab_root
[FTP] 220 LabFTP Server Ready
[FTP] USER anonymous -> 331 Password required
[FTP] PASS anon@lab -> 230 Login successful
[FTP] PASV -> 227 Entering Passive Mode (192,168,0,116,195,84)
[FTP] LIST -> 150 Directory listing
[FTP] STOR uploads/shell.py -> 226 Transfer complete (Web shell uploaded)
[FTP] DELE uploads/shell.py -> 250 File deleted successfully (Forensic trace removed)
\end{logbox}

\subsection{Master Empirical Verification Matrix}
Table~\ref{tab:master_matrix} summarizes the definitive empirical verification runs across our
laboratory testbed, cross-referencing target systems, operational environments, and observed results
without relying on raw file labels:

\begin{table}[!ht]
\centering
\footnotesize
\renewcommand{\arraystretch}{1.08}
\caption{Consolidated Master Matrix of Empirical Verification Test Runs}
\label{tab:master_matrix}
\begin{tabularx}{\textwidth}{@{}L{3.8cm} L{2.4cm} L{2.6cm} X@{}}
\toprule
\hdrrow \thh{Empirical Evaluation Scenario} & \thh{Target Platform} & \thh{Defense State} & \thh{Key Attribution \& Observed Outcome}\\
\midrule
Full Attack Chain \& Active Defense & Windows 11 Desktop & Active Blocker On & Windows 62\%; 3/4 exploits succeeded; severed mid-attack on Run 2.\\
\addlinespace[2pt]
Dynamic OS-Aware Exploitation & Win 11 + Kali 6.x & Passive / Off & Dynamic payload swap: \texttt{win.ini}/\texttt{whoami} vs.\ \texttt{/etc/passwd}/\texttt{id}.\\
\addlinespace[2pt]
Closed-Port RST \& ICMP Evaluation & Kali Linux 6.x VM & Passive / Off & Linux 76\% (Closed-port RST alone) / 67\% (ICMP Code 9 fallback).\\
\addlinespace[2pt]
100-Port Comprehensive Range Scan & Windows 11 Host & Passive / Off & Batch 1: Win 71\% (SYN/ACK); Batch 2: Win 69\% (Zero open ports RST).\\
\addlinespace[2pt]
Darwin / BSD Mobile Stack Telemetry & Apple iPad (iPadOS) & Passive / Off & macOS/BSD 48\% (Darwin option order: WScale before Timestamp).\\
\addlinespace[2pt]
Android Linux Mobile Stack Telemetry & Android 14 Phone & Passive / Off & Linux 58\% (SACK before TS, RFC 6864 all-zero IP-ID with DF set).\\
\addlinespace[2pt]
OpenSSH Weak-Credential Compromise & Windows OpenSSH & Passive / Off & Windows 75\%; cracked \texttt{test:test123}; loud desktop payloads deployed.\\
\addlinespace[2pt]
Windows Firewall Dynamic Quarantine & Windows 11 Desktop & Windows Firewall & Tripped at 8 SYNs; canary port 81 trapped; auto-released after 60s TTL.\\
\addlinespace[2pt]
Linux Netfilter Dynamic Mitigation & Kali Linux VM & Linux Netfilter & Dynamic Netfilter DROP rule; drop counter progressed 12 $\to$ 16 pkts.\\
\bottomrule
\end{tabularx}
\end{table}

\subsection{Cross-Platform Network Stack Signature Comparison}
Table~\ref{tab:os_comparison} details the empirical behavioral metrics that allowed the engine to
differentiate Microsoft Windows from Unix-like operating systems:

\begin{table}[!ht]
\centering
\footnotesize
\renewcommand{\arraystretch}{1.12}
\caption{Cross-Platform Stack Behavioral Signatures Observed in Empirical Testbed}
\label{tab:os_comparison}
\begin{tabularx}{\textwidth}{@{}L{5.0cm} >{\RaggedRight\arraybackslash}X >{\RaggedRight\arraybackslash}X@{}}
\toprule
\hdrrow \thh{Diagnostic Metric} & \thh{Windows 11 Desktop} & \thh{Kali Linux / Android}\\
\midrule
Initial IPv4 TTL & 128 & 64\\
TCP Window Size & 65535 or 8192 & 29200, 65160, or 0 (closed RST)\\
IP-ID Sequence Archetype & \texttt{STRICT\_INCREMENTAL} ($\Delta \approx 1\dots 3$) & \texttt{ALL\_ZERO} ($\mathrm{ID} = 0$, RFC 6864 DF=1)\\
TCP Timestamp Clock Rate & Calibrated $\approx 1000$\,Hz (1\,ms tick) & $\approx 250$\,Hz or 1000\,Hz\\
RFC 3168 ECN Negotiation & Ignored (\texttt{0x12} returned) & Supported (\texttt{0x52} ECE reflected)\\
Invalid Flag Probe (SYN+FIN) & Emits TCP RST on open port & Silently dropped without reply\\
ICMP Non-Standard Code 9 & Normalized to Code 0 & Reflected unaltered as Code 9\\
OS-Aware Traversal Target & \texttt{windows\textbackslash win.ini} & \texttt{/etc/passwd}\\
RCE Command Execution & \texttt{whoami} $\to$ \texttt{homed01\textbackslash turjo} & \texttt{id} $\to$ \texttt{uid=0(root)}\\
\bottomrule
\end{tabularx}
\end{table}
\FloatBarrier
